\paragraph{window-\texorpdfstring{$6$}{6} 群商实现的稳定子分层与 Lee--Yang 振幅常数的双二次阴影}
在 window-\(6\) 的纤维退化直方图、自由作用与线性核障碍已经锁定之后，任何试图把纤维分解解释为有限群轨道分解的方案都必须携带多层稳定子；另一方面，Lee--Yang 振幅常数 \(B\) 的真正六次性并不止于 \(\QQ(B)/\QQ\) 的次数陈述，其分裂域还显现出一条与三次分辨子域不同的额外二次阴影。

\begin{theorem}[window-\texorpdfstring{$6$}{6} 的任何群商实现都强制三层稳定子分层]\label{thm:xi-window6-group-quotient-isotropy-strata}
设 \(\Omega_6:=\{0,1\}^6\)，并令
\[
F_6:=\Fold^{\mathrm{bin}}_6:\Omega_6\to X_6.
\]
若存在某个有限群 \(G\) 在 \(\Omega_6\) 上作用，使其轨道分解恰等于 \(F_6\) 的纤维分解，则
\[
\boxed{\;12\mid |G|.\;}
\]
并且轨道稳定子必分成三种不同的阶层：对每个 \(\omega\in\Omega_6\)，有
\[
|\mathrm{Orb}_G(\omega)|\in\{2,3,4\}
\quad\Longleftrightarrow\quad
|\mathrm{Stab}_G(\omega)|\in\left\{\frac{|G|}{2},\frac{|G|}{3},\frac{|G|}{4}\right\}.
\]
更具体地，大小分别为 \(2,3,4\) 的轨道个数恰为
\[
8,\qquad 4,\qquad 9.
\]
因此任何这样的群商实现都不可能是自由作用，而只能是带有三层稳定子分层的离散 orbifold 型商。
\leanverified{paper\_xi\_window6\_group\_quotient\_isotropy\_strata}
\end{theorem}
\begin{proof}
由推论 \ref{cor:foldbin6-degeneracy-multivalued-nonfree}，\(F_6\) 的纤维划分既不可能来自自由有限群作用，也不可能来自阿贝尔线性陪集划分。另一方面，由推论 \ref{cor:fold-bin-m6-fiber-hist}，
\[
|\{w\in X_6:\ |F_6^{-1}(w)|=2\}|=8,\qquad
|\{w\in X_6:\ |F_6^{-1}(w)|=3\}|=4,\qquad
|\{w\in X_6:\ |F_6^{-1}(w)|=4\}|=9.
\]
若群 \(G\) 的轨道分解恰与纤维分解一致，则轨道大小只能取 \(2,3,4\)，且对应轨道数正是上述 \(8,4,9\)。

对任意 \(\omega\in\Omega_6\)，轨道--稳定子定理给出
\[
|\mathrm{Orb}_G(\omega)|=\frac{|G|}{|\mathrm{Stab}_G(\omega)|}.
\]
因此 \(2,3,4\) 都必须整除 \(|G|\)，遂有
\[
\mathrm{lcm}(2,3,4)=12\mid |G|.
\]
并且当轨道大小分别为 \(2,3,4\) 时，稳定子阶只能分别为
\[
\frac{|G|}{2},\qquad \frac{|G|}{3},\qquad \frac{|G|}{4}.
\]
这三者彼此不同，故点稳定子至少分成三层；又因为全部轨道只出现这三种大小，所以稳定子阶也恰只出现这三层。证毕。
\end{proof}

\begin{theorem}[Lee--Yang 振幅常数 \texorpdfstring{$B$}{B} 的六次性与双二次阴影]\label{thm:xi-leyang-amplitude-B-sextic-biquadratic-shadow}
设
\[
f_B(X):=162X^6+1593X^4+1998X^2+1184.
\]
则
\[
\boxed{\;[\QQ(B):\QQ]=6.\;}
\]
并且 \(f_B\) 在 \(\QQ\) 上不可约，正是 \(B\) 的六次最小多项式；其判别式满足
\[
\boxed{\;
\Disc(f_B)=
-2^{16}\cdot 3^{22}\cdot 11^4\cdot 37^3\cdot 109^4,\;}
\]
故平方自由部分为
\[
\boxed{\;-37.\;}
\]
若 \(N\) 为 \(f_B\) 的分裂域，则
\[
\QQ(\sqrt{-37})\subset N.
\]
另一方面，\(N\) 同时包含 \(B^2\) 的三次最小多项式的分裂域，因而也包含其唯一二次分辨子域
\[
\QQ(\sqrt{-111}).
\]
于是
\[
\boxed{\;
\QQ(\sqrt{-37},\sqrt{-111})\subset N,
\qquad
12\mid [N:\QQ].\;}
\]
\end{theorem}
\begin{proof}
由定理 \ref{thm:xi-leyang-amplitude-B-true-quadratic-cover}，
\[
[\QQ(B):\QQ]=6.
\]
又因 \(B^2\) 满足三次方程
\[
g(T):=162T^3+1593T^2+1998T+1184,
\]
故 \(g(B^2)=0\)，亦即 \(f_B(B)=g(B^2)=0\)。由于 \(f_B\) 的次数也是 \(6\)，它必为 \(B\) 的最小多项式，从而在 \(\QQ\) 上不可约。

设 \(r_1,r_2,r_3\) 为 \(g\) 的根，则
\[
f_B(X)=162\prod_{j=1}^{3}(X^2-r_j),
\]
其六个根为 \(\pm \sqrt{r_1},\pm \sqrt{r_2},\pm \sqrt{r_3}\)。于是
\[
\Disc(f_B)
=
162^{10}\Bigl(\prod_{j=1}^{3}4r_j\Bigr)
\prod_{1\le i<j\le 3}(r_i-r_j)^4.
\]
另一方面，
\[
\Disc(g)=162^4\prod_{1\le i<j\le 3}(r_i-r_j)^2,
\qquad
\prod_{j=1}^{3}r_j=-\frac{1184}{162}.
\]
故可化为
\[
\Disc(f_B)
=
-64\cdot 162\cdot 1184\cdot \Disc(g)^2.
\]
再由推论 \ref{cor:xi-leyang-u-b-square-common-quadratic-resolvent}，
\[
\Disc(g)=
-2^2\cdot 3^9\cdot 11^2\cdot 37\cdot 109^2.
\]
代入即得
\[
\Disc(f_B)=
-2^{16}\cdot 3^{22}\cdot 11^4\cdot 37^3\cdot 109^4.
\]
因此其平方自由部分为 \(-37\)。分裂域 \(N\) 因而包含
\[
\QQ(\sqrt{\Disc(f_B)})=\QQ(\sqrt{-37}).
\]

再者，\(N\) 含有 \(f_B\) 的全部根 \(\pm\sqrt{r_j}\)，平方后即得到 \(g\) 的全部根 \(r_j\)，故 \(g\) 的分裂域 \(L_b\) 满足
\[
L_b\subset N.
\]
由推论 \ref{cor:xi-leyang-u-b-square-common-quadratic-resolvent}，\(L_b\) 的唯一二次子域是
\[
\QQ(\sqrt{-111}).
\]
于是 \(N\) 同时包含 \(\QQ(\sqrt{-37})\) 与 \(\QQ(\sqrt{-111})\)。由于这两个虚二次域不同，遂得
\[
\QQ(\sqrt{-37},\sqrt{-111})\subset N.
\]
又因 \([L_b:\QQ]=6\)，且 \(\QQ(\sqrt{-37})\not\subset L_b\)（因为 \(L_b\) 的唯一二次子域为 \(\QQ(\sqrt{-111})\)），故
\[
[L_b(\sqrt{-37}):\QQ]=12.
\]
最后注意 \(L_b(\sqrt{-37})\subset N\)，于是
\[
12\mid [N:\QQ].
\]
证毕。
\end{proof}

\endinput
